\chapter{研究内容}
\section{总体思路：以实际损失为共同监督目标}
研究围绕一个共同目标组织不同信息，而不是为价格、公告和图分别定义互不相同的风险。预测日$t$的证券$i$具有过去价格历史、市场与行业环境，以及已公开事件和持股关系。模型输出未来五个市场日累计收益的四档分布，由此得到发生亏损、达到$5\%$损失和达到$10\%$损失的嵌套概率。不同来源对同一未来结果提供条件信息，它们不承担彼此替代的标签。

价格信息既是持续更新的底座，也是理解公告和关系的环境。同一经营预告在上涨、普跌或高波动时期，可能对应不同后续损失概率；共同机构下的多只证券，也可能主要受到同一行业冲击。本研究将个股、市场和行业分别编码成64、32和32维历史状态，拼接得到128维共同价格上下文。事件注意力以该上下文为query，关系消息读取邻居截至当日已知的价格与事件状态，三个专家均接收共同环境。

\reportfigure{共同价格上下文与可靠性约束的四档专家混合}{fig:architecture}{\FigArchitecture}{批准研究设计。各维度、掩码和概率接口为拟实施合同，图不表示正式风险模型已经训练；蓝、橙、绿框均有直接文字及边框编码。}

图\ref{fig:architecture}体现了两个设计选择。其一，价格专家不只是输出一个等待小幅修补的logit，而是一个完整的四档专家；事件和关系专家也能表达强损失情形。其二，辅助专家的最大权重由资料可靠性限制，而不是由它们自身给出的风险高低决定。强事件能否主导，取决于可用、质量、年龄和学习门，不由任意固定logit帽限制。这一结构仍属于混合专家思路，创新假设集中在时点资料约束的具体构造和检验。

\section{研究问题与证据安排}
首先研究共同价格状态对实际损失的预测能力。个股历史波动与共同市场环境都有合理经济依据，但GRU不必然超过滚动收益和下行统计。研究将保留透明Logit和强树，比较个股价格、价格加市场行业以及完整价格上下文模型，分析共同因素是否尤其有助于普跌日期。去市场或去行业控制同时关闭相应显式派生特征，不允许删除一个分支后继续访问同一信息。

其次研究公告语义和持股关系的增量。公告方向是经营语义，不是未来收益标签；持股关系是公开的信息关联，不是已识别的因果传播。原生字段、预测结构和正文表示分层比较，可以区分直接数值、生成抽取与连续语义的作用。真实关系、去关系和保持度数及行业组成的重连，可以检查图增量是否只是网络容量或行业构成。各项比较使用相同预测总体和相同已知结果评价集合。

第三研究可靠性约束相对普通门控、拼接及标准概率MoE的净作用。约束可以给出缺失回退和年龄相关的原始概率偏移界，但这些性质不是泛化保证。正式比较需要观察AP、NLL、Brier、校准与预警漏报，并对自然年龄、低质量和低共享邻居分组分析。若约束只在一种覆盖群体有效，研究应解释条件而不是仅给跨年平均增幅。

\reporttable{研究问题与主要证据对象}{tab:questions}{
\begin{tabularx}{\linewidth}{p{2.8cm}YY}\toprule
研究问题&主要比较&需要观察的证据\\\midrule
共同价格环境的作用&透明统计、强树、个股与市场行业GRU&逐年$p_5$排序、概率质量、共同下跌漏报\\
公告语义的增量&原生字段、预测结构、正文、去事件&独立语义误差、覆盖与同信息下游差值\\
共享关系的增量&去关系、真实图、度数/行业组成保持重连&共享邻居、去自身边界与概率增量\\
可靠性约束的作用&拼接、普通门控、标准MoE与本方法&同目标校准、缺失回退、年龄界及压力\\
总体与结果分离&认证总体、观察候选及未知原因&全总体覆盖、监督子集与风险广度界限\\\bottomrule
\end{tabularx}}

\section{数据产品：预测索引与结果索引分离}
数据工作从原始资料的经济对象和历史时点开始。不同工作簿可能含同义表头，也可能使用不同公司码；原始名单可能混合总股东、流通股东及补充记录。标准化应保留源文件指纹、字段与行位置，用确定性映射形成候选产品。它的目的不是让所有缺失和冲突都变成“可用”，而是使合法连接与不确定状态可以明确识别。

预测索引以时点上市A股总体为主，记录证券、预测日、当时身份和交易状态。过去输入通过市场日序列索引、事件公开版本缓存和完整名单版本流按日查询。未来结果另表计算，左连接到预测索引后形成四档监督或未知原因。图\ref{fig:products}中的分离让模型输入不必通过一张已筛除未来缺失的宽表生成，也防止邻居节点隐含获得未来标签。

\reportfigure{时点总体、过去输入与离线结果的产品关系}{fig:products}{\FigProducts}{数据接口设计。预测成员由当时状态建立，结果可空；未来结果不参与历史输入查询和总体筛选。}

来源认证不足时，工程执行可以产生“观察总体/完整窗口候选面板”。这个名称反映能够从现有源行确认的范围，不声称覆盖所有当时上市证券。补齐股票主表、每日交易估值状态和终端价值之后，才判断哪些候选可以进入认证面板。候选数量需要实际运行生成，不能从原残差标签规模推算，更不能把原样本ID直接改名沿用。

\section{公告与关系表示的具体研究内容}
公告工作复用已锁定的1500条语料、800/200/500划分以及200试标和300双标。语义目标保持方向、明确原因和逐字证据，标注者不读取未来风险标签。Qwen完成适配和独立语义评价后冻结，风险任务使用其预测结构和正文向量，而非测试金标。正文池化排除系统、包装、padding和生成答案；缓存键将原文、模型、分词器、模板和正文范围绑定。

事件表示将合法正文与合法JSON分开。解析失败会影响结构字段，但未必使正文隐藏表示无效；证据可在正文找到，也不意味着它完整支持了语义答案。不同状态应进入质量记录和覆盖报告，不能合并成单一成功率。价格上下文query对最近可用事件进行选择，学习在当前环境下哪些事件更相关，最终事件专家仍预测同一四档结果。

关系工作以接受的完整公开名单为更新单位，区分主体与历史证券对应。新完整版本替换旧名单，Top10退出只说明离开披露范围，不等于清仓。比例和边龄形成非负权重，机构聚合邻居已知价格与事件状态，再扣除目标自身贡献。没有其他正权共享邻居时关系不可用。这样的接口才能支持“事件与价格状态沿关系聚合”的表述，只有传播价格向量的旧原型不应被描述为事件传播。

\section{训练、校准与输出之间的关系}
训练采用最终混合四档概率的NLL，并使用匹配的模态dropout和共同正则化。价格专家在辅助缺失样本上接受同一主目标的学习，事件与图专家也由混合目标联合优化。旧分支辅助BCE系数不进入主方案；若研究其影响，作为另行登记、同目标匹配的敏感性实验。这样，同编码器比较主要改变融合机制，而不是同时改变监督预算。

验证年1--8月用于选参和早停，9--12月用于最终四档温度校准及预警阈值，两个阶段按结果结束日剔除跨界窗口。温度作用于最终混合分布，不假定各专家先校准后混合自然可靠。全部档位使用同一组混合权重，校准后仍由四档求和获得嵌套概率。原始年龄界和非线性校准后的概率质量分别检查。

\reporttable{主要输入与输出的张量合同}{tab:tensors}{
\begin{tabularx}{\linewidth}{p{3.2cm}p{3.15cm}Y}\toprule
对象&形状或维度&含义\\\midrule
个股历史&$N\times60\times F_s$&实际收益及仅过去派生量，逐字段掩码\\
市场历史&$D\times60\times F_M$&市场原始收益、FF3及过去状态，每日共享\\
行业历史&$C\times60\times F_I$&时点已知门类环境，类别日期共享\\
共同价格上下文&$N\times128$&个股64、市场32、行业32拼接\\
正文与事件状态&$N\times K\times D_Q$&$K\le8$，正文与预测结构分别判定有效\\
二部边及邻居状态&$2\times E$；节点状态向量&比例、边龄、历史身份与价格+事件状态\\
专家与混合分布&三个$N\times4$及一个$N\times4$&非负归一，所有档位使用相同权重\\
结果与状态&$N$&四档监督可空，未知原因与监督掩码另存\\\bottomrule
\end{tabularx}}

表\ref{tab:tensors}的$F_s,F_M,F_I$由冻结字段顺序决定，不能沿用旧五字段序列或合成示例默认值来代表新装配。原始输入值、掩码、来源可用日和监督状态各自需要检查。最终原型保留全总体的$p_{loss},p_5,p_{10}$、专家分布和混合权重，按绝对阈值预警；每天Top10\%只是次要审核容量指标。市场和行业平均$p_5,p_{10}$用于估计损失成员占比，结合未知结果上下界评价共同下跌覆盖，不扩写为指数事件概率。
